\documentclass[12pt]{article}
\usepackage[utf8]{inputenc}
\usepackage[T1]{fontenc}
\usepackage[margin=1in]{geometry}
\usepackage{setspace}
\usepackage{xcolor}
\usepackage[colorlinks=true, urlcolor=blue!60!black, linkcolor=blue!60!black]{hyperref}
\setlength{\parindent}{0em}
\setlength{\parskip}{0.45em}
\renewcommand{\baselinestretch}{1.0}
\begin{document}
\begin{center}
{\Large \textit{Palko v.\ Connecticut}, 302 U.S.\ 319 (1937)\par}
\vspace{0.3em}
{\normalsize Case Brief --- One-Page Edition\par}
\end{center}

\textbf{Memory Jogger.} Only rights implicit in ordered liberty bind the States; Connecticut could retry a murder defendant (overruled by \textit{Benton}, 1969) (\href{https://supreme.justia.com/cases/federal/us/302/319/}{Justia}).

\textbf{Facts.} Frank Palka killed two Bridgeport officers in 1935. Convicted of second-degree murder, he was retried after the State appealed legal errors under an 1886 statute, convicted of first-degree murder, and sentenced to death.

\textbf{Procedural History.} The Connecticut Supreme Court of Errors ordered the retrial and affirmed the death sentence.

\textbf{Votes.} Affirmed, 8--1: Cardozo for the Court; Butler dissented without opinion.

\textbf{Holding.} The Fourteenth Amendment absorbs only rights ``of the very essence of a scheme of ordered liberty''; a retrial to correct legal error is not intolerably shocking.

\textbf{Reasoning.} Rejecting ``the tyranny of labels,'' Cardozo noted that States could drop grand juries, juries, and the self-incrimination privilege, while free speech, the ``matrix'' of other freedoms, and counsel in capital cases were absorbed. Connecticut sought only ``a trial free from the corrosion of substantial legal error,'' not to wear Palko down after an error-free trial.

\textbf{Examples.} \emph{Same-Side:} charging by information. \emph{Opposite-Side:} a prior restraint on newspapers. \emph{Fence-Sitter:} non-unanimous juries (later barred in \textit{Ramos}).

\textbf{Critique.} Total-incorporation originalists (Black, Amar, Thomas) fault its reliance on due process and its judicial picking of rights; civil libertarians found its ``civilized system'' benchmark stingy toward defendants.

\textbf{Key Quotation.} ``In these and other situations, immunities that are valid as against the federal government by force of the specific pledges of particular amendments have been found to be implicit in the concept of ordered liberty, and thus, through the Fourteenth Amendment, become valid as against the states.'' --- Justice Cardozo, 302 U.S.\ at 324--25.

\end{document}
